\documentclass[10pt,a4paper]{amsart}
\usepackage[margin=19mm]{geometry}
\usepackage{amsmath,amssymb,mathtools}
\usepackage[scr=rsfs,cal=euler]{mathalfa}
\usepackage{xfrac}
\usepackage[unicode,hidelinks,bookmarks=false]{hyperref}
\newcommand{\ie}{i.e.\ }
\newcommand{\cf}{cf.\ }
\newcommand{\resp}{respectively\ }
\newcommand{\Subsection}{\subsection}
\providecommand{\nobreakdash}{\nobreak\hbox{-}}
\setlength{\emergencystretch}{2em}
\multlinegap0pt
\usepackage{enumerate}
\usepackage{amssymb}
\usepackage{xcolor}
\newtheorem{lemma}{Lemma}
\newtheorem{thm}{Theorem}
\newcommand{\Hom}{{\mathrm{Hom}}}
\newcommand{\SO}{{\mathrm{SO}}}
\title{Howe duality for the dual pair $SL_2(\mathbb R) \times F_{4,1}$: a ping-pong of $K$-types}
\author{Gordan Savin}
\date{Source: arXiv:2508.12534v1 (CC BY 4.0)}
\begin{document}
\maketitle

\noindent\emph{Source and licence.} Howe duality for the dual pair $SL_2(\mathbb R) \times F_{4,1}$: a ping-pong of $K$-types. Version arXiv:2508.12534v1 (CC BY 4.0), distributed by arXiv under the Creative Commons Attribution 4.0 International licence, http://creativecommons.org/licenses/by/4.0/, as recorded in the arXiv licence metadata for this exact version and shown on the abstract page https://arxiv.org/abs/2508.12534v1. Full licence notice: \texttt{licenses/arxiv-2508.12534-CC-BY-4.0.txt}.\par\smallskip

\noindent\textit{Editorial scope.} Counted unit: the article's thm thm (source0.tex lines 342--346). The article's complete original proof of it is delivered with it (source0.tex lines 347--365, one \texttt{proof} environment of 19 lines). No mathematics is written, repaired or re-derived: the statement and the proof are the article's own lines, in the article's own notation, and no retained line is substituted or re-flowed.  Retained uncounted as the source-local context the proof needs: lines 299--304.  Nothing was omitted from any retained span, so the package carries no omission record.  No retained line contains a citation, so the wrapper carries no bibliography and no bibliography entry.  No retained line contains a figure environment, an image-inclusion command or drawing code, so no image is delivered and the package declares no asset dependency.  Editorial formatting only, in addition: an AMS \texttt{amsart} preamble that restates the article's own declarations and macros used by the retained lines, namely the environments \texttt{lemma}, \texttt{thm} on separate counters, together with the standard packages those lines need.  The retained environments print in this document as Lemma 1, Theorem 1 in order of appearance, whereas the article prints the same items with the numbering of its own full document.  The complete native source of the article as distributed in the arXiv e-print for this exact version is bundled unchanged as \texttt{sources/183416/source0.tex}.
\begin{lemma} \label{L:1} 
 Let $\sigma$ be a finite length $(\mathfrak{sl}_2,\SO(2))$-module. Then 
\[ 
\Hom_K(\Theta(\sigma), \tau(m,n)) \cong \Hom_{\mathfrak{sl}_2} (\Theta(\tau(m,n)), \sigma) \cong \Hom_{\SO(2)}( (2m+4) , \sigma). 
\] 
\end{lemma} 

\begin{thm} Let $\sigma$ be an irreducible $(\mathfrak{sl}_2,\SO(2))$-module. Assume that $\sigma$ contains the type $(2m+4)$, for $m\geq 0$, and no smaller types $2n+4$, with $n\geq 0$. 
 (This condition is empty if $m=0$.) Then $\Theta(\sigma)$ has a unique irreducible quotient. It contains the type $\tau(m,0)$ with multiplicity one, and no types $\tau(n,0)$ with $n<m$. 
Conversely, if $\pi$ is an irreducible $(\mathfrak g,K)$-module containing the type $\tau(m,0)$, and 
no smaller such types, then $\Theta(\pi)$ has unique irreducible quotient. It contains the type $(2m+4)$, and no smaller types $2n+4$, with $n\geq 0$. 
\end{thm} 

\begin{proof}  Assume $\pi$ is a quotient of $\Theta(\sigma)$. We do not assume that $\pi$ is irreducible. By Lemma \ref{L:1}, we have the following sequence of inclusions: 
\[ 
\Hom_K(\pi, \tau(m,0)) \subseteq \Hom_K(\Theta(\sigma), \tau(m,0))  \cong \Hom_{\mathfrak{sl}_2} (\Theta(\tau(m,0)), \sigma) \cong \Hom_{\SO(2)} ((2m+4) , \sigma). 
\] 
Observe that we can ran this sequence with any $2n+4$ in place of $2m+4$. If $n<m$, by the assumption, the last space is trivial, hence $\tau(n,0)$ is not a type of 
$\pi$. We shall use this in a moment. 
Since $\sigma$ is a quotient of $\Theta(\pi)$ we have the second sequence of inclusions (note that we are starting with the space of the same dimension as 
$\Hom_{\SO(2)} ((2m+4) , \sigma)$: 
\[ 
\Hom_{\SO(2)}(\sigma, (2m+4)) \subseteq \Hom_{\SO(2)}(\Theta(\pi), (2m+4))  \cong \Hom_{\mathfrak g} (\Theta(2m+4), \pi) \subseteq  \Hom_{K} (F_m , \pi). 
\] 
Since, $\pi$ has no type $\tau(n,0)$ with $n<m$, we ended with  $\Hom_{K} (\tau(m,0) , \pi)$, which has  the same dimension as $\Hom_K(\pi, \tau(m,0))$. 
Thus all inclusions in the two sequences  are isomorphisms, and thus all spaces are one-dimensional, since $\Hom_{\SO(2)} ((2m+4) , \sigma)$ is one-dimensional. 
However, we did not assume that $\pi$ is irreducible. If $\pi=\pi_1\oplus \pi_2$ then, 
\[ 
1+1 = \dim \Hom_{K}(\pi_1,   \tau(m,0)) + \dim\Hom_{K}(\pi_2,  \tau(m,0)) = \dim  \Hom_{K}(\pi, \tau(m,0))=1
\] 
a contradiction. Thus $\Theta(\sigma)$ has a unique irreducible quotient. It contains $\tau(m,0)$, with multiplicity one.  The other direction is proved in the same way... 
\end{proof} 

\end{document}
